\chapter{球面、雙曲空間與\(k\)的意義}

我們先做一點簡單的數學。任意維度的球面都是特殊的幾何：它們具有均勻性，也就是曲率為常數。先看正曲率，再看均勻負曲率，接著再說明它們與愛因斯坦方程的關係。

球形宇宙是指\emph{空間}是球形的，而不是時空。另一種可能性是具有均勻負曲率。把這兩種情形並列比較，才能看清廣義相對論的作用。

\section{球面與均勻幾何}

一維球面就是圓。描述它最簡單的方式，是將它嵌入二維輔助空間，並寫成
\begin{equation}
x^2+y^2=R^2.
\end{equation}
單位一維球面則是
\begin{equation}
x^2+y^2=1.
\end{equation}

二維球面也用同樣的方式描述，只是現在要用三個輔助座標：
\begin{equation}
x^2+y^2+z^2=R^2,
\end{equation}
單位球面的情形為
\begin{equation}
x^2+y^2+z^2=1.
\end{equation}

這些座標只是輔助工具。我們關心的幾何並不是周圍的\(x,y,z\)空間，而是球面本身。生活在一維球面上的生物生活在圓上，無法接觸圓的內部或外部。二維球面也是如此。球面內部的填充部分和球面外部，都只是輔助圖像的一部分。

同樣的構造也能得到三維球面。再引入一個座標，例如\(w\)，並寫成
\begin{equation}
x^2+y^2+z^2+w^2=R^2,
\end{equation}
單位球面的情形為
\begin{equation}
x^2+y^2+z^2+w^2=1.
\end{equation}
三維球面較難想像，因為輔助工具現在是四維的，但數學並沒有困難多少。

標準記號是用\(\Omega_n\)表示\(n\)維球面。因此，\(\Omega_1\)是一維球面，\(\Omega_2\)是二維球面，\(\Omega_3\)是三維球面。

\section{球面上的度量}

更實用的描述方式是內在描述，也就是用定義在球面本身的座標來描述球面。

描述單位一維球面只需要一個座標\(\phi\)，其範圍為\(0\leq \phi \leq 2\pi\)。它的度量為
\begin{equation}
ds^2=d\phi^2.
\end{equation}
也可以寫成
\begin{equation}
d\Omega_1^2=d\phi^2.
\end{equation}

單位二維球面可以視為一系列的一維球面。角度\(\phi\)仍沿著每個圓繞行，而另一個角度\(\theta\)則從球面的一端延伸到另一端。這個角度的範圍是\(0\)到\(\pi\)。沿\(\theta\)方向的位移對度量的貢獻是\(d\theta^2\)。而沿\(\phi\)方向繞行時，所在圓的半徑為\(\sin\theta\)。因此
\begin{equation}
d\Omega_2^2=d\theta^2+\sin^2\theta\,d\phi^2
\end{equation}
或寫成
\begin{equation}
d\Omega_2^2=d\theta^2+\sin^2\theta\,d\Omega_1^2.
\end{equation}

三維球面也以同樣的方式構造，現在可視為一系列的二維球面。引入另一個座標\(\alpha\)，則
\begin{equation}
d\Omega_3^2=d\alpha^2+\sin^2\alpha\,d\Omega_2^2.
\end{equation}
展開後為
\begin{equation}
d\Omega_3^2=d\alpha^2+\sin^2\alpha\left(d\theta^2+\sin^2\theta\,d\phi^2\right).
\end{equation}
此處，\(\alpha\)從一端的零尺寸延伸到另一端，因此\(0\leq \alpha \leq \pi\)。\editorialnote{00：17：24時，音訊聽起來短暫像是說了``零到二{\proseunicodefont π}''；但依照整個構造，以及明確出現的\(\sin^2\alpha\)因子，新座標必須從這一系列低維球面由零尺寸開始的一端，延伸至另一端。}

一般的遞迴形式為
\begin{equation}
d\Omega_n^2=d\alpha^2+\sin^2\alpha\,d\Omega_{n-1}^2.
\end{equation}

採用這種觀點後，就可以不再理會嵌入圖像。一維球面由一個座標描述，二維球面由兩個座標描述，三維球面由三個座標描述，以此類推。這種幾何本身就是完整的。

這種幾何之所以特殊，是因為它具有齊性。若將\(\sin^2\theta\)換成另一個在兩端仍為零的函數，就能改變度量，得到變形後的閉曲面，例如橢球面。但這樣的曲面就不再均勻。球面具有齊性，因為每一點都與其他各點相同。

若要描述半徑為\(A\)而非單位半徑的球面，只要將單位球面的度量乘上\(A^2\)：
\begin{equation}
ds^2=A^2\,d\Omega_3^2.
\end{equation}
必須取平方，因為度量衡量的是距離的平方。

順帶一提，零維球面可以由構造向前推一步得到：
\begin{equation}
x^2=1,
\end{equation}
其解為兩個點\(x=\pm 1\)。因此，單位零維球面就是兩個點。

\section{球形 FRW 宇宙}

現在關心的不只是空間，而是時空。相關的圖像是一個半徑隨時間變化的三維球面。

有一點必須始終謹記：時間是不同的。不要把膨脹中的宇宙想成高維球面，並把時間當作另一個嵌入座標。

有了這個認識，時空間隔為
\begin{equation}
d\tau^2=dt^2-a(t)^2\,d\Omega_3^2.
\end{equation}
這是空間為球形的宇宙所對應的 FRW 度量。相較於靜態球面，唯一的新成分是半徑變成了時間的函數\(a(t)\)。

將我們選在某個極點上，會讓座標較方便。那個位置並沒有任何特殊之處，這只是座標約定。如此一來，\(\alpha\)可以視為單位球面上從我們所在位置向外的徑向座標。它不是以公里為單位的徑向距離，而是單位球面上的無因次角度量，只是在幾何中扮演徑向距離的角色。

除星系、恆星等局部團塊之外，這就是空間為球形的 FRW 宇宙的完整幾何。它的空間曲率為正，符號約定寫作
\begin{equation}
k=+1.
\end{equation}
這本身並不是``加一''，而是表示空間曲率為正。

\section{從球面向外觀看}

從球面上的某一點向外看，這樣的球面會呈現什麼樣子？這就要用到球極投影。

取一個在南極與無限平面相切的二維球面。在球面上選一點，從北極穿過該點畫一條直線，並延長至與平面相交。如此一來，球面上的每一點都會映射到平面上的一點。南極會映射到平面的原點。北極則是例外：從不同方向接近北極時，它的像會沿著不同方向延伸至無窮遠。因此，我們將北極視為無窮遠點。

若在球面不同位置放上大小相同的``大陸''，它們在平面上的像會在接近北極時逐漸變大。因此，若觀察者位於南極，遠方物體所張的角會比在平坦空間中更大。在正曲率空間中，已知大小的物體相對於歐幾里得幾何看起來會過大。

這就是基本的觀測概念：根據距離、物體大小與角大小之間的關係推斷曲率。星系本身並不是測量極遠距離的理想標準尺，但推理方式正是如此。

\begin{classroomqa}
\lecturetimestamp{00:45:54}
\audiencequestion{除了角大小之外，還需要知道距離嗎？}
\lecturerresponse{需要。光知道角度還不夠，也必須知道距離。}
\end{classroomqa}

\begin{classroomqa}
\lecturetimestamp{00:46:18}
\audiencequestion{球面上的圓投影後仍會是圓嗎？}
\lecturerresponse{會。圓會投影成圓，這是球極投影的一項特殊性質。}
\end{classroomqa}

\begin{classroomqa}
\lecturetimestamp{00:47:02}
\audiencequestion{大陸經過投影後會維持原有形狀嗎？}
\lecturerresponse{這裡有兩個不同的說法。一個是圓會投影成圓；另一個較弱的說法是，足夠小的形狀會保留其局部形狀，因為角度不變。這種映射是保角的：它保留角度，但不保留相對大小。}
\end{classroomqa}

\section{均勻負曲率}

均勻負曲率的情形，會從度量開始介紹。

球面的負曲率類比是雙曲空間。依照前面相同的遞迴構造，將\(\sin\alpha\)換成\(\sinh\alpha\)，便可得到它的度量：
\begin{equation}
d\mathcal H_n^2=d\alpha^2+\sinh^2\alpha\,d\Omega_{n-1}^2.
\end{equation}
低維的角向部分仍是單位球面度量\(d\Omega_{n-1}^2\)。

\begin{classroomqa}
\lecturetimestamp{00:55:51}
\audiencequestion{度量中的低維微分項仍然是球面的那一項嗎？}
\lecturerresponse{是。低維部分就是一般的單位球面度量。改變的是它前面的\(\sinh^2\alpha\)因子。}
\end{classroomqa}

雙曲正弦為
\begin{equation}
\sinh\alpha=\frac{e^\alpha-e^{-\alpha}}{2}.
\end{equation}
負號很重要；講座中一位學生隨即指出並更正了這個符號。

座標範圍也有一項關鍵差異。在球面上，新座標從一端延伸至另一端，並在\(\pi\)處結束。在這裡，
\begin{equation}
0\leq \alpha<\infty.
\end{equation}
一般正弦先增大，然後又回到零。雙曲正弦在\(\alpha=0\)附近起初也有相似的行為，但接著會呈指數增長。

這表示，沿著這種幾何向外移動時，周圍的低維球面會愈來愈大。在這種幾何中，半徑增加時，天空不會先變大再縮小，而是持續擴大。

在三個空間維度中，相應的 FRW 形式為
\begin{equation}
d\tau^2=dt^2-a(t)^2\,d\mathcal H_3^2,
\end{equation}
其曲率符號標記為
\begin{equation}
k=-1.
\end{equation}

\section{圓盤圖像及其意義}

負曲率可以透過埃舍爾的圓盤圖像直觀呈現，尤其是\emph{《極限圓 IV》}。它和球極投影的類比很有用，但兩者的表現恰好相反。

球面的投影看起來無限延伸，但實際幾何是有限的；負曲率的圓盤圖像看起來有限，但實際幾何是無限的。大小相同的``大陸''或``星系''愈接近邊界，看起來就愈小，但在幾何上，它們的大小完全相同。若以經過多少個等大區塊來計量距離，邊界就位於無窮遠處。

因此，在負曲率空間中，已知大小的物體相對於平坦空間看起來太小，所張的角也太小。

這種幾何仍然具有均勻性。雖然圖像中外側區域看起來很小，但若將觀察者移到其他位置，他看到的內在幾何仍會相同。

\begin{classroomqa}
\lecturetimestamp{01:06:42}
\audiencequestion{如果觀察者位於其中任一個圓內，也會看到類似的圖樣嗎？}
\lecturerresponse{會，完全相同。空間具有均勻性。}
\end{classroomqa}

\begin{classroomqa}
\lecturetimestamp{01:07:03}
\audiencequestion{在球面情形中，遠處物體看起來太大；在這裡則看起來太小嗎？}
\lecturerresponse{是。在正曲率空間中，已知大小的物體相對於平坦空間所張的角太大；在負曲率空間中，所張的角太小。}
\end{classroomqa}

\begin{classroomqa}
\lecturetimestamp{01:08:09}
\audiencequestion{這種類似球極投影的投影，是否能讓天文學家從觀測結果重建幾何？}
\lecturerresponse{是。天文學家實際上就像測量師，測量天空中的三角形。就目前要說明的幾何概念而言，可以先不考慮時間變化，專注於角度和距離。}
\end{classroomqa}

還有一個有用的局部圖像：可以把一小塊負曲率空間畫成普通三維空間中的馬鞍面。但只能這樣表示局部區域；整個雙曲平面無法完整嵌入普通三維歐幾里得空間，成為一個形狀良好的曲面。

\section{平坦情形}

介於正曲率與負曲率情形之間的是平坦空間。可以把它看成一種極限，不過只要將它視為一般的平坦歐幾里得空間即可。

它的 FRW 形式為
\begin{equation}
d\tau^2=dt^2-a(t)^2\left(dx^2+dy^2+dz^2\right),
\end{equation}
其中
\begin{equation}
k=0.
\end{equation}

這就是三種均勻 FRW 幾何：
\begin{align}
k=+1 &: \text{positive spatial curvature},\\
k=0 &: \text{flat space},\\
k=-1 &: \text{negative spatial curvature}.
\end{align}

\begin{classroomqa}
\lecturetimestamp{01:15:07}
\audiencequestion{若將幾何投影到平面上並追蹤一個訊號，投影中光速會隨位置改變嗎？}
\lecturerresponse{會，在投影圖像中會如此。在球面投影中，愈往外看，光似乎移動得愈快；在負曲率圓盤中，愈接近邊緣，光似乎移動得愈慢。這種變化只出現在圖像中，並不屬於底層物理。}
\end{classroomqa}

\begin{classroomqa}
\lecturetimestamp{01:17:05}
\audiencequestion{人們真的曾用這種方式觀察遠方物體來判定空間是否平坦嗎？}
\lecturerresponse{有，但實際使用的是角度資訊。判定空間是否平坦，並不是靠測量光在投影圖像中看起來移動得多快。}
\end{classroomqa}

\section{回到宇宙學方程}

休息之後，接下來要說明先前研究過的宇宙學方程與愛因斯坦方程之間的關係。此時黑板上仍留有平坦空間度量的一部分。

\begin{figure}[htbp]
\centering
\includegraphics[width=0.72\textwidth]{lecture_05_figure_02.png}
\caption{平坦 FRW 度量片段。}
\lectureframe{01:18:18}
\end{figure}

首先簡短回顧牛頓力學。把我們放在中心，觀察一個依照哈勃定律遠離我們的星系。根據牛頓的球對稱分布定理，只有該星系距離所對應半徑的球體內部質量，會對它所受的力有所貢獻。於是有三種可能：星系的運動速度高於逃逸速度、低於逃逸速度，或恰好等於逃逸速度。

接著回顧，對一般非相對論物質而言，密度與質量除以\(a^3\)成正比；對輻射而言，密度則與\(1/a^4\)成正比；而真空能的密度不會隨球體大小改變。

先前的牛頓力學處理中，宇宙學方程包含一個密度項；若不是恰好處於逃逸速度的情形，還會多出一個正比於\(1/a^2\)的常數項。先前的討論大多將此項設為零。從經驗上看，此項很小，但在概念上沒有理由將它省略。

\section{愛因斯坦方程與\(00\)分量}

接下來看愛因斯坦場方程。左側是由時空曲率構成的幾何，右側則是作為引力場源頭的能量動量張量。按照講座採用的歸一化，方程寫成
\begin{equation}
R^{\mu\nu}-\frac{1}{2}g^{\mu\nu}R=4\pi G\,T^{\mu\nu}.
\end{equation}
\editorialnote{黑板上可見的方程中，因子\(\tfrac12\)無法辨認，但在01:33:36附近，講者口頭更正時補上了這個因子。此處的方程採用更正後的形式，並保留講座中的\(4\pi G\)歸一化。}

此處不需要用到十個獨立分量中的每一個。到目前為止，我們已經作了許多假設---包括均勻性、FRW 形式，以及只有一個隨時間變化的尺度因子---因此只剩下一個任意函數，也就是\(a(t)\)。自然的選擇是時間—時間分量，因為\(T^{00}\)就是能量密度\(\rho\)。

\begin{figure}[htbp]
\centering
\includegraphics[width=0.62\textwidth]{lecture_05_figure_05.png}
\caption{愛因斯坦場方程及其\(00\)分量。}
\lectureframe{01:30:54}
\end{figure}

因此，焦點是
\begin{equation}
R^{00}-\frac{1}{2}g^{00}R=4\pi G\,\rho.
\end{equation}

右側是能量密度，左側是幾何。

\begin{classroomqa}
\lecturetimestamp{01:39:09}
\audiencequestion{能量動量張量的其他分量呢？}
\lecturerresponse{一般而言，它們不為零。空間—空間分量可以表示壓力。在物質主導的情形中，壓力為零；在輻射主導的情形中，壓力是能量密度的三分之一；對真空能，則寫作\(P=w\rho\)，其中\(w=-1\)。}
\end{classroomqa}

\begin{classroomqa}
\lecturetimestamp{01:40:24}
\audiencequestion{所以，如果只有物質，\(T^{\mu\nu}\)的其他分量會消失嗎？}
\lecturerresponse{對無壓物質而言，會，壓力分量為零。}
\end{classroomqa}

\section{愛因斯坦計算的結果}

講者沒有在黑板上完整推導曲率。重點在於將 FRW 度量代入愛因斯坦張量後會出現什麼結果。

其中一類項只來自尺度因子隨時間變化這件事。即使空間是平坦的，隨時間變化的尺度因子仍會產生時空曲率。這項貢獻正比於
\begin{equation}
\left(\frac{\dot a}{a}\right)^2.
\end{equation}

第二類項來自空間本身的曲率。即使尺度因子固定不隨時間變化，彎曲空間仍會對時空曲率有所貢獻。對於簡單且均勻彎曲的幾何，曲率與曲率半徑的平方成反比，因此會產生一個正比於
\begin{equation}
\frac{k}{a^2},
\end{equation}
的項；球面情形為正號，雙曲情形為負號，平坦空間則為零。

\begin{classroomqa}
\lecturetimestamp{01:36:53}
\audiencequestion{是否還應該有像曲率半徑那樣的因子？}
\lecturerresponse{那已經包含在\(1/a^2\)因子裡了。對這些簡單而均勻彎曲的幾何而言，曲率就是曲率半徑平方的倒數。}
\end{classroomqa}

推導過程中即席處理了係數，也部分參照先前的牛頓力學結果進行核對。最後得到\footnote{這個弗里德曼方程形式，整合了講者口頭討論並自行修正的係數，以及在01:32:35--01:36:53附近與先前牛頓力學結果所作的比較。}
\begin{equation}
\left(\frac{\dot a}{a}\right)^2+\frac{k}{a^2}=\frac{8\pi G}{3}\rho,
\end{equation}
或等價地寫成
\begin{equation}
\left(\frac{\dot a}{a}\right)^2=\frac{8\pi G}{3}\rho-\frac{k}{a^2}.
\end{equation}

這個形式將三種 FRW 幾何與先前討論的宇宙動力學連結起來。

\begin{classroomqa}
\lecturetimestamp{01:37:52}
\audiencequestion{如果重新定義\(a\)，\(k\)可以取任意值嗎？}
\lecturerresponse{在平坦情形下，\(a\)沒有唯一自然的尺度，所以可以重新縮放，但\(k\)仍為零。在球面與負曲率情形下，則存在自然的曲率尺度。若令\(a(t)\)等於這個自然的曲率半徑，\(k\)就只是\(+1\)或\(-1\)。}
\end{classroomqa}

\begin{classroomqa}
\lecturetimestamp{01:41:00}
\audiencequestion{在最後的方程中，\(k\)等於\(+1\)、\(-1\)還是\(0\)？}
\lecturerresponse{是。球面情形為\(+1\)，雙曲情形為\(-1\)，平坦空間則為\(0\)。}
\end{classroomqa}

牛頓力學的處理沒有誤導我們；它得到了相同的基本方程。廣義相對論補上的是詮釋。從牛頓力學的觀點來看，額外的常數用來區分運動速度高於逃逸速度、低於逃逸速度，以及恰好等於逃逸速度的臨界情形。從愛因斯坦的觀點來看，同一個符號則代表空間本身的幾何。

\begin{classroomqa}
\lecturetimestamp{01:42:46}
\audiencequestion{還沒有納入宇宙學常數時，這套關係是如何成立的？}
\lecturerresponse{在引入宇宙學常數之前，逃逸速度的運動情形與空間幾何之間有緊密的對應關係。若\(k=+1\)，空間曲率為正，運動速度低於逃逸速度，因此宇宙可以先膨脹再收縮。若\(k=-1\)，空間曲率為負，運動速度高於逃逸速度，因此宇宙會持續膨脹。若\(k=0\)，則是臨界的刀鋒情形。}
\end{classroomqa}

因此，同一個符號\(k\)可以從兩種角度解讀：一是運動的動力學分類，二是空間曲率的符號。
